../cictikz/src/cictikz/data/tex/cictikz_preamble.tex

% An OFDM link, drawn deliberately as the single carrier link with
% different inputs.
%
% The blocks after the transform are identical: there is still one
% carrier, one transmitter, one antenna. What changed is that the thing
% being modulated is now a set of sub-carriers described in frequency
% space, and an inverse transform turns that description into the I and
% Q the transmitter wants. The receiver runs the transform the other
% way. The vertical dots stand for the sub-carriers between the first
% and the last.

\begin{tikzpicture}[thick]
  %%%%%%%%%%%%%%%%%%%%%%%%%%%%%%%%%%%%%%%%%%%%%%%%%%%%%%%%%%%%%%%%%%%%%%
%% Shared pieces for the l04_afe signal flow graphs, l4_first_order and
%% l4_biquad.  The biquad is meant to read as the first-order graph with a
%% second integrator added, so the summing node and the 1/s block have to be
%% the same size and shape in both.  They are the same size in the original
%% artwork too.
%%
%% Names are prefixed and letters only, and computed coordinates are braced:
%% TikZ scans a coordinate for its closing parenthesis, so a '(' inside an
%% expression ends it early.
%%%%%%%%%%%%%%%%%%%%%%%%%%%%%%%%%%%%%%%%%%%%%%%%%%%%%%%%%%%%%%%%%%%%%%

\newcommand{\sfgR}{0.25}    % summing node radius
\newcommand{\sfgBw}{1.0}    % half width of the 1/s block
\newcommand{\sfgBh}{0.65}   % half height of the 1/s block

% A summing node at (#1,#2).
\newcommand{\sfgSum}[2]{%
  \draw (#1,#2) circle (\sfgR);
  \draw ({#1-0.17},#2) -- ({#1+0.17},#2);
  \draw (#1,{#2-0.17}) -- (#1,{#2+0.17});
}

% A block centred at (#1,#2) carrying #3.
\newcommand{\sfgBox}[3]{%
  \draw ({#1-\sfgBw},{#2-\sfgBh}) rectangle ({#1+\sfgBw},{#2+\sfgBh});
  \node at (#1,#2) {#3};
}

\newcommand{\sfgDot}[2]{\fill (#1,#2) circle (0.07);}


  \def\yt{1.25}
  \def\yb{-1.25}

  % ---- transmitter ----
  \draw[->] (-2.4,\yt) -- (0,\yt);
  \draw[->] (-2.4,\yb) -- (0,\yb);
  \node[armygreen, anchor=east, scale=0.8] at (-2.5,\yt)
    {$a_1(t)e^{j(\omega_1 t + \phi_1(t))}$};
  \node[armygreen, anchor=east, scale=0.8] at (-2.5,\yb)
    {$a_N(t)e^{j(\omega_N t + \phi_N(t))}$};
  \foreach \y in {-0.4,0,0.4} { \fill (-1.2,\y) circle (0.055); }

  \sfgTall{0}{-1.8}{1.3}{1.8}{IFFT}
  \draw[->] (1.3,\yt) -- (2.7,\yt) node[midway, above] {$I$};
  \draw[->] (1.3,\yb) -- (2.7,\yb) node[midway, above] {$Q$};
  \sfgTall{2.7}{-1.8}{4.0}{1.8}{TX}
  \draw (4.0,0) -- (5.0,0);
  \sfgAntenna{5.0}{0}

  % ---- receiver ----
  \sfgAntenna{7.4}{0}
  \draw (7.4,0) -- (8.4,0);
  \sfgTall{8.4}{-1.8}{9.7}{1.8}{RX}
  \draw[->] (9.7,\yt) -- (11.1,\yt) node[midway, above] {$I$};
  \draw[->] (9.7,\yb) -- (11.1,\yb) node[midway, above] {$Q$};
  \sfgTall{11.1}{-1.8}{12.4}{1.8}{FFT}
  \draw[->] (12.4,\yt) -- (13.4,\yt);
  \draw[->] (12.4,\yb) -- (13.4,\yb);
  \foreach \y in {-0.4,0,0.4} { \fill (12.9,\y) circle (0.055); }
  \node[armygreen, anchor=west, scale=0.8] at (13.5,\yt)
    {$a_1(t)e^{j(\omega_1 t + \phi_1(t))}$};
  \node[armygreen, anchor=west, scale=0.8] at (13.5,\yb)
    {$a_N(t)e^{j(\omega_N t + \phi_N(t))}$};
\end{tikzpicture}
\end{document}
